\chapter{我为什么教书}\label{s:finale}

我开始在多伦多大学做志愿者时，
有些学生问我为什么要免费教书。
这是我的回答：

\begin{quote}

我在你们这个年纪时，
以为大学的存在是为了教人如何学习。
后来，
在研究生院，
我以为大学是关于做研究、创造新知识的。
如今我四十多岁了，
却意识到我们真正在教你们的，
是如何接管这个世界，
因为不管你们愿不愿意，你们都不得不这么做。

我的父母七十多岁了。
他们不再掌管这个世界；
制定法律、
在医院里做生死决定的，是我这个年纪的人。
尽管很吓人，
但\emph{我们}才是成年人。

二十年后，
我们会走向退休，而\emph{你们}将掌权。
在你们十九岁时，这听起来可能很遥远，
但喘三口气，它就过去了。
这就是为什么我们给你们出那些没法抄去年笔记的题。
这就是为什么我们把你放进那样的处境里：
你必须弄清眼下什么该做、
什么可以留到以后、
什么你可以干脆不管。
因为如果你们现在不学会做这些事，
到了不得不做的时候，你们就没准备好。

\end{quote}

这些全是真的，
但并不是故事的全部。
我并不希望别人把世界变得更好，好让我能安逸地退休。
我希望他们去做，是因为这是我们这个时代最伟大的冒险。
一百五十年前，
大多数社会还在实行奴隶制。
一百年前，
我祖母在加拿大\hreffoot{https://en.wikipedia.org/wiki/The\_Famous\_Five\_(Canada)}{在法律上还算不上一个人}。
在我出生的那一年，
世界上大多数人在极权统治下受苦，
法官还在判人去接受电击疗法来「治愈」同性恋。
这个世界仍然有很多毛病，
但看看我们比祖辈多了多少选择。
看看我们因为终于认真对待「金律」，
能多知道多少东西、能成为多少种人、能享受多少乐趣。

我今天比那时更不乐观了。
气候变化、
物种大灭绝、
监视资本主义、
一个世纪以来未曾见过的大规模不平等、
种族主义民族主义的卷土重来：
我这一代人眼看着这一切发生，却耸耸肩。
我们懦弱、懈怠和贪婪的账单，要等我女儿长大才到期，
但它\emph{一定}会来，
而等到它来的时候，这些问题就不会有容易的解决办法了
（甚至可能根本没有办法）。

所以这就是我今天教书的原因：
我愤怒。
我愤怒，是因为你的性别、你的肤色、你父母的财富和关系，
不该比你的聪明、诚实或勤奋更重要。
我愤怒，是因为我们把互联网变成了污水池。
我愤怒，是因为纳粹再次上街，
而亿万富翁在地球融化的时候玩火箭。
我愤怒，
所以我教书，
因为只有当我们教人如何把世界变得更好，世界才会变好。

在他 1947 年的文章「我为什么写作」里，
\hreffoot{http://www.resort.com/~prime8/Orwell/whywrite.html}{George Orwell}写道：
\index{Orwell, George}

\begin{quote}

  在和平年代，我可能写的是辞藻华丽或纯粹描述性的书，
  而且可能几乎察觉不到自己的政治立场。
  但现实是，我被迫成了一种小册子作者{\ldots}
  自 1936 年以来我写的每一行严肃作品，
  无论直接还是间接，
  都是反对极权主义的{\ldots}
  在我们这样的时代，
  认为一个人可以不去写这类题材，在我看来是无稽之谈。
  每个人都在以这样或那样的面貌写它们。
  问题只在于一个人站在哪一边。

\end{quote}

\noindent
把「写作」换成「教书」，你就得到了我做这件事的原因。

\vspace{\baselineskip}

\noindent
谢谢你读到这里——希望我们有一天能一起教书。
在那之前：

\begin{center}

从哪里开始都行。\\
有什么就用什么。\\
能帮谁就帮谁。

\end{center}
